\clearpage\thispagestyle{otFront}\phantomsection\label{mode-emploi}\pdfbookmark[0]{Mode d’emploi}{bm-mode-emploi}\otPageTitre{Mode d’emploi}{Des repères pour travailler seul}\fontsize{11}{15}\selectfont

\begin{tabularx}{\linewidth}{@{}lY@{}}\toprule
Une pastille & Découverte : vocabulaire, lecture, calcul direct.\\[3mm]
Deux pastilles & Entraînement : choisir une méthode, organiser des étapes.\\[3mm]
Trois pastilles & Problèmes : relier les notions, justifier et interpréter.\\\bottomrule
\end{tabularx}\par\medskip
Ces indications décrivent les tâches ; elles ne classent pas les élèves.\par\medskip
\begin{exemple}[Lire une correction]Pour $5x=18$, on cherche le nombre
qui, multiplié par 5, donne 18. Donc $x=18:5=3{,}6$.
Le contrôle $5\times3{,}6=18$ confirme le résultat.
Refaire cette explication aide davantage que recopier seulement la réponse.\end{exemple}
\begin{attention}Lire les unités, les hypothèses et les mots « au moins »,
« strictement » ou « total ». Une figure ne donne pas à elle seule les
propriétés d’un objet. Une probabilité décrit un modèle de hasard,
pas une prédiction certaine du prochain résultat.\end{attention}
Chercher avant d’ouvrir le corrigé. Après lecture, fermer le livre et refaire
le passage décisif sur une feuille vierge. Les liens permettent de passer
de l’énoncé à sa correction puis de revenir.

\clearpage
